% !TEX root = ../main.tex
\section{Experiments}
\label{sec:exp}

\subsection{Setup}
\label{sec:setup}
We examine LTF and DiffusionDrive (front camera and ego status), DiffusionDriveV2 (camera and LiDAR),
SimLingo (InternVL2-1B), AutoVLA (Qwen2.5-VL-3B, three cameras $\times$ four frames) and Alpamayo-1.5 (four cameras, a diffusion head over the vision--language model cache), all with released weights. The three BEV planners are trained on NAVSIM (nuPlan) data, as is AutoVLA's released checkpoint. SimLingo is trained in CARLA and Alpamayo-1.5 on NVIDIA's PhysicalAI-AV recordings. None is fine-tuned here, so every policy is read on a corpus and sensor rig it was not trained on. For the BEV planners the nuScenes front image is centre-cropped to the 4:1 aspect of NAVSIM's stitched front view with the sky removed and resized to the network's input size, the logged speed and acceleration fill the status vector, and DiffusionDriveV2 receives the LiDAR sweep. The VLA policies receive the nuScenes cameras mapped approximately onto the views they expect. DiffusionDrive's diffusion seed is fixed so that paired queries share their noise.

The primary corpus is nuScenes~\cite{caesar2020nuscenes} keyframes: corridor pedestrians (lateral
clearance to the ego corridor $\le1$\,m) at $5$--$12$\,m, their off-road counterparts, and $30$
approach sequences at $5.8$--$14.6$\,m with logged pedestrian futures. The near-pedestrian units come from $88$ Boston (left-hand-drive) and $148$ Singapore (right-hand-drive) front-only frames. Alpamayo-1.5 covers $170$ of them, needing $1.6$\,s of history and $6.4$\,s of future
within the scene. \Cref{sec:axes} repeats $\mathcal{E}_R$ and $\mathcal{E}_N$ on a second, right-hand-drive corpus, the \NumSGscenes{} NAVSIM test-split~\cite{dauner2024navsim} Singapore scenes with a pedestrian near the ego corridor and a complete logged pedestrian future, used for the transfer and decomposition results and for the NAVSIM columns of \cref{tab:report}. \Cref{tab:report} is measured on the nuScenes frames alone.
